%%%%%%%%%%%%%%%%%%%%%%%%%%%%%%%%%%%%%%%%%%%%%%%%%%%%%%%%%%%%%%%
%
% Welcome to Overleaf --- just edit your LaTeX on the left,
% and we'll compile it for you on the right. If you open the
% 'Share' menu, you can invite other users to edit at the same
% time. See www.overleaf.com/learn for more info. Enjoy!
%
%%%%%%%%%%%%%%%%%%%%%%%%%%%%%%%%%%%%%%%%%%%%%%%%%%%%%%%%%%%%%%%


% Inbuilt themes in beamer
\documentclass{beamer}

\usepackage[ngerman]{babel}
\usepackage{siunitx}

\usepackage{biblatex}
\usepackage{filecontents}
\usepackage{subcaption}
\usepackage{appendixnumberbeamer}

\bibliography{ref_pres.bib}



% Theme choice:
\usetheme{Berlin}
\usecolortheme{beaver}
\setbeamertemplate{page number in head/foot}{\insertframenumber/\inserttotalframenumber} 

% Title page details: 

\title{Counting wheat heads on minimaly annotated data using AttentionMIL} 
\subtitle{PMasterarbeit}
\author{Erik Marten Wegener}
\date{06.06.2026}

\definecolor{dred}{rgb}{0.8, 0.0, 0.0}

\setbeamertemplate{navigation symbols}{}

\setbeamercolor{itemize item}{fg=dred}
\setbeamercolor{itemize subitem}{fg=dred}
\setbeamercolor{itemize subsubitem}{fg=dred}

\setbeamertemplate{enumerate items}[default]
\setbeamercolor{enumerate item}{fg=dred}
\setbeamercolor{enumerate subitem}{fg=dred}
\setbeamercolor*{enumerate subsubitem}{fg=dred}

\setbeamertemplate{itemize item}[square]
\setbeamertemplate{itemize subitem}[circle]
\setbeamertemplate{itemize subsubitem}[triangle]

\setbeamercolor{section number projected}{bg=gray}

\setbeamercolor{block title}{
	use=frametitle,fg=frametitle.fg,
	use=frametitle right,bg=frametitle right.bg!60}
\setbeamercolor{block body}{use=frametitle,bg=frametitle.bg}

\setbeamerfont{caption}{size=\scriptsize}

\setbeamercolor{caption name}{fg=dred}
\begin{document}
	
	% Title page frame
	
	

	\begin{frame}[plain]
		\titlepage
	\end{frame}
	
	
	% Remove logo from the next slides
	\logo{}
	
	
	% Outline frame
	
	\begin{frame}{Gliederung}
		\addtocounter{framenumber}{-1}
		\tableofcontents
	\end{frame}

	\section{Motivation}
		\begin{frame}{Motivation}
			\begin{figure}[H]
				\centering
				\includegraphics[width=\textwidth]{Bilder/Drone_example.JPG}
			\end{figure}
		\end{frame}

		\begin{abstract}
			
		\end{abstract}
	
	% Lists frame
	\section{The idea}
	% Present the idea from Ilse et al with the attention weights for the different numbers
	% Lead to the idea of couning high attention patches in order to make an estimate
	% Excoures to MIL, Attention and AttentionMIL
	%	MIL: Multiple Instance Learning is just the solving of a problem by not looking at each instance seperately but rather looking at a bag of instances and making the decsion
	%	if there is a leath one positive instance in the bag or not. This is a weakly supervised learning method.
	%	Attention: Attention is a method to make a model focus on the most important parts. The attention weights are learned 
	%   during training and can be used to visualize what the model is focusing on.
	%	AttentionMIL: AttentionMIL is a combination of MIL and Attention. The model is trained to focus on the most important instances 
	%   in the bag and make the decsion based on those instances.
	% show results on MNIST 

	% \begin{frame}{Performance-Vergleich zu bekannten Modellen(1)}
	% 	\begin{figure}[H] 
	% 		\begin{subfigure}[b]{0.66\linewidth}
	% 			\centering
	% 			\includegraphics[width=5cm]{} 
	% 			\label{fig:DLPRB-Performance} 
	% 			\vspace{4ex}
	% 		\end{subfigure}%% 
	% 		\begin{subfigure}[b]{0.33\linewidth}
	% 			\centering
	% 			\includegraphics[width=3cm]{} 
	% 			\label{} 
	% 			\vspace{4ex}
	% 		\end{subfigure}  
	% 		\caption{Muster}  
	% 	\end{figure}
	% \end{frame}

	
	\begin{frame}{Literatur}
		\tiny
		\printbibliography
	\end{frame}
	
\end{document}